% ══════════════════════════════════════════════════════════════════════════════
%                         فصل پنجم
%                 کمونیسم: از مارکس تا مائو
% ══════════════════════════════════════════════════════════════════════════════

\chapter{کمونیسم: از مارکس تا مائو}
\label{ch:communism}

\begin{flushright}
\textit{«یک شبح در اروپا پرسه می‌زند—شبح کمونیسم.»}

— کارل مارکس و فردریش انگلس، مانیفست کمونیست (۱۸۴۸)
\end{flushright}

% ══════════════════════════════════════════════════════════════════════════════
%                         مقدمه فصل
% ══════════════════════════════════════════════════════════════════════════════

\lettrine[lines=3, lraise=0.1, nindent=0.5em]{\textcolor{LeftRedDark}{ک}}{مونیسم} 
در قرن بیستم از یک نظریه فلسفی به نیرویی تبدیل شد که بر یک‌سوم جمعیت جهان حکومت کرد. از انقلاب اکتبر ۱۹۱۷ تا فروپاشی شوروی در ۱۹۹۱، کمونیسم هم الهام‌بخش میلیون‌ها انسان بود و هم توجیه‌گر برخی از بزرگ‌ترین جنایات تاریخ. در این فصل، سیر تحول کمونیسم را از نظریه به عمل، و از مارکس تا مائو دنبال می‌کنیم.

% ══════════════════════════════════════════════════════════════════════════════
\section{کمونیسم به مثابه هدف نهایی}
% ══════════════════════════════════════════════════════════════════════════════

\subsection{تفاوت سوسیالیسم و کمونیسم در نظریه مارکسی}

در نوشته‌های مارکس، «سوسیالیسم» و «کمونیسم» گاه به جای هم به کار می‌روند. اما تمایزی نیز وجود دارد:

\begin{defbox}[تعریف: سوسیالیسم و کمونیسم در نظریه مارکس]
\textbf{سوسیالیسم (مرحله پایین):}
\begin{itemize}[nosep]
    \item مالکیت جمعی ابزار تولید برقرار شده
    \item اما هنوز دولت وجود دارد (دیکتاتوری پرولتاریا)
    \item توزیع بر اساس «کار»: هرکس به اندازه کارش دریافت می‌کند
    \item هنوز کمیابی و نابرابری‌هایی وجود دارد
\end{itemize}

\textbf{کمونیسم (مرحله بالا):}
\begin{itemize}[nosep]
    \item دولت «زوال» یافته—جامعه خودگردان
    \item طبقات کاملاً محو شده‌اند
    \item فراوانی مادی تحقق یافته
    \item اصل توزیع: «از هرکس به اندازه توانش، به هرکس به اندازه نیازش»
    \item کار، دیگر اجبار نیست بلکه لذت است
\end{itemize}
\end{defbox}

\begin{quotebox}[کارل مارکس، نقد برنامه گوتا (۱۸۷۵)]
\textit{«در مرحله بالاتر جامعه کمونیستی، پس از آنکه انقیاد برده‌وار افراد به تقسیم کار از میان رفت... کار نه فقط وسیله زندگی، بلکه خود نخستین نیاز زندگی شد... آنگاه جامعه می‌تواند بر پرچم خود بنویسد: از هرکس به اندازه توانش، به هرکس به اندازه نیازش!»}
\end{quotebox}

\subsection{کمونیسم: آرمان‌شهر یا علم؟}

\begin{debatebox}[آیا کمونیسم قابل تحقق است؟]
\textbf{مارکسیست‌ها:}
\begin{itemize}[nosep]
    \item کمونیسم نتیجه منطقی تکامل تاریخی است
    \item با رشد نیروهای مولده، فراوانی ممکن می‌شود
    \item با پایان استثمار، دولت بی‌معنا می‌شود
\end{itemize}

\textbf{منتقدان:}
\begin{itemize}[nosep]
    \item انسان ذاتاً خودخواه است—بدون انگیزه، کار نمی‌کند
    \item کمیابی همیشه وجود دارد (حتی اگر مادی نباشد، زمان و توجه کمیاب‌اند)
    \item بدون دولت، چگونه اختلافات حل می‌شوند؟
    \item تجربه قرن بیستم: دولت کمونیستی نه‌تنها زوال نیافت، بلکه قوی‌تر شد
\end{itemize}
\end{debatebox}

% ══════════════════════════════════════════════════════════════════════════════
\section{لنینیسم: نظریه و عمل}
% ══════════════════════════════════════════════════════════════════════════════

\person{ولادیمیر ایلیچ لنین} مارکسیسم را از نظریه به انقلاب پیروزمند تبدیل کرد. اما در این مسیر، آن را نیز تغییر داد.

\begin{personbox}[ولادیمیر ایلیچ لنین (۱۸۷۰-۱۹۲۴)]

\textbf{نام اصلی:} ولادیمیر ایلیچ اولیانوف

\textbf{ملیت:} روسی

\textbf{آثار کلیدی:} \book{چه باید کرد؟} (۱۹۰۲)، \book{امپریالیسم، بالاترین مرحله سرمایه‌داری} (۱۹۱۶)، \book{دولت و انقلاب} (۱۹۱۷)

\mediumspace

\textbf{نوآوری‌های لنین:}
\begin{enumerate}[nosep]
    \item \textbf{حزب پیشاهنگ:} انقلاب نیاز به حزب متمرکز و منضبط از انقلابیون حرفه‌ای دارد
    \item \textbf{امپریالیسم:} سرمایه‌داری با استثمار جهان سوم، بحران‌هایش را به تعویق انداخته
    \item \textbf{ضعیف‌ترین حلقه:} انقلاب لزوماً در پیشرفته‌ترین کشورها رخ نمی‌دهد
    \item \textbf{دموکراسی‌متمرکز:} بحث آزاد، اما اجرای متحد پس از تصمیم
\end{enumerate}

\mediumspace

\textbf{نقل‌قول:} \textit{«بدون تئوری انقلابی، عمل انقلابی ممکن نیست.»}
\end{personbox}

\subsection{حزب پیشاهنگ}

\begin{defbox}[تعریف: حزب پیشاهنگ لنینی]
بر خلاف احزاب سوسیال‌دموکرات که توده‌ای بودند، لنین خواستار:
\begin{itemize}[nosep]
    \item حزب کوچک از \textbf{انقلابیون حرفه‌ای}
    \item انضباط آهنین و تمرکز
    \item عضویت محدود و سخت‌گیرانه
    \item رهبری روشنفکران که «آگاهی طبقاتی» را به کارگران می‌برند
\end{itemize}
\textbf{دلیل:} کارگران خودبه‌خود فقط به «آگاهی صنفی» (خواست دستمزد بیشتر) می‌رسند، نه آگاهی انقلابی.
\end{defbox}

\subsection{انقلاب اکتبر ۱۹۱۷}

\begin{figure}[htbp]
\centering
\begin{tikzpicture}[scale=0.85]

% عنوان
\node[font=\large\bfseries, text=DarkGray] at (0, 7.5) {\fa{تایم‌لاین: انقلاب روسیه (۱۹۱۷)}};

% خط زمان
\draw[line width=4pt, MediumGray] (-7, 5.5) -- (7, 5.5);

% رویدادها
% انقلاب فوریه
\fill[OrangeAccent] (-6, 5.5) circle (10pt);
\node[above, text=OrangeAccent, font=\small\bfseries] at (-6, 6) {\fa{فوریه}};
\node[below, text width=2.2cm, align=center, text=DarkGray, font=\scriptsize] at (-6, 4.8) {
    \fa{سقوط تزار}\\
    \fa{حکومت موقت}
};

% بازگشت لنین
\fill[LeftRed] (-3.5, 5.5) circle (10pt);
\node[above, text=LeftRed, font=\small\bfseries] at (-3.5, 6) {\fa{آوریل}};
\node[below, text width=2.2cm, align=center, text=DarkGray, font=\scriptsize] at (-3.5, 4.8) {
    \fa{بازگشت لنین}\\
    \fa{تزهای آوریل}
};

% قیام ژوئیه
\fill[MediumGray] (-1, 5.5) circle (10pt);
\node[above, text=MediumGray, font=\small\bfseries] at (-1, 6) {\fa{ژوئیه}};
\node[below, text width=2.2cm, align=center, text=DarkGray, font=\scriptsize] at (-1, 4.8) {
    \fa{قیام ناموفق}\\
    \fa{لنین مخفی شد}
};

% کودتای کورنیلوف
\fill[RightBlue] (1.5, 5.5) circle (10pt);
\node[above, text=RightBlue, font=\small\bfseries] at (1.5, 6) {\fa{اوت}};
\node[below, text width=2.2cm, align=center, text=DarkGray, font=\scriptsize] at (1.5, 4.8) {
    \fa{کودتای کورنیلوف}\\
    \fa{شکست راست}
};

% انقلاب اکتبر
\fill[LeftRedDark] (4.5, 5.5) circle (12pt);
\node[above, text=LeftRedDark, font=\small\bfseries] at (4.5, 6.1) {\fa{۲۵ اکتبر}};
\node[below, text width=2.5cm, align=center, text=DarkGray, font=\scriptsize] at (4.5, 4.7) {
    \fa{انقلاب اکتبر}\\
    \fa{بلشویک‌ها قدرت}\\
    \fa{را گرفتند}
};

% جنگ داخلی
\fill[LeftRedDark] (6.5, 5.5) circle (10pt);
\node[above, text=LeftRedDark, font=\small\bfseries] at (6.5, 6) {\fa{۱۸-۲۱}};
\node[below, text width=2cm, align=center, text=DarkGray, font=\scriptsize] at (6.5, 4.8) {
    \fa{جنگ داخلی}\\
    \fa{پیروزی سرخ‌ها}
};

% شعار
\node[text=LeftRed, font=\bfseries] at (0, 3.5) {\fa{«نان، صلح، زمین! تمام قدرت به شوراها!»}};

\end{tikzpicture}
\caption{سیر انقلاب روسیه در سال ۱۹۱۷}
\label{fig:russian-revolution}
\end{figure}

\begin{historybox}[انقلاب اکتبر: انقلاب یا کودتا؟]
در ۲۵ اکتبر ۱۹۱۷ (۷ نوامبر تقویم جدید)، بلشویک‌ها کاخ زمستانی را تصرف کردند.

\textbf{موافقان «انقلاب»:}
\begin{itemize}[nosep]
    \item شوراها (سویت‌ها) از بلشویک‌ها حمایت می‌کردند
    \item حکومت موقت مشروعیت نداشت
    \item توده‌ها خواستار صلح و زمین بودند
\end{itemize}

\textbf{موافقان «کودتا»:}
\begin{itemize}[nosep]
    \item بلشویک‌ها اقلیت بودند (حتی در شوراها)
    \item مجلس مؤسسان منتخب را منحل کردند
    \item با خشونت قدرت را حفظ کردند
\end{itemize}
\end{historybox}

% ══════════════════════════════════════════════════════════════════════════════
\section{استالینیسم}
% ══════════════════════════════════════════════════════════════════════════════

پس از مرگ لنین (۱۹۲۴)، \person{یوسف استالین} در مبارزه قدرت پیروز شد و تا ۱۹۵۳ بر شوروی حکومت کرد. دوران او، کمونیسم را برای همیشه تغییر داد.

\begin{personbox}[یوسف ویساریونوویچ استالین (۱۸۷۸-۱۹۵۳)]

\textbf{نام اصلی:} یوسف جوگاشویلی

\textbf{ملیت:} گرجی

\textbf{دوران حکومت:} ۱۹۲۴-۱۹۵۳

\mediumspace

\textbf{ویژگی‌های استالینیسم:}
\begin{enumerate}[nosep]
    \item \textbf{سوسیالیسم در یک کشور:} می‌توان در شوروی سوسیالیسم ساخت، بدون انقلاب جهانی
    \item \textbf{صنعتی‌سازی اجباری:} برنامه‌های پنج‌ساله، تبدیل شوروی به قدرت صنعتی
    \item \textbf{اشتراکی‌سازی کشاورزی:} کولاک‌ها (دهقانان مرفه) نابود شدند
    \item \textbf{پاک‌سازی‌ها:} میلیون‌ها نفر اعدام، زندانی یا تبعید شدند
    \item \textbf{پرستش شخصیت:} استالین «پدر ملت‌ها» شد
\end{enumerate}

\mediumspace

\textbf{میراث:} صنعتی‌سازی سریع، پیروزی در جنگ جهانی دوم، اما به بهای جان میلیون‌ها انسان
\end{personbox}

\begin{table}[htbp]
\centering
\caption{هزینه انسانی استالینیسم (تخمین‌های مورخان)}
\label{tab:stalin-victims}
\begin{tabular}{r|r}
\toprule
\textbf{رویداد} & \textbf{تعداد قربانیان (تخمین)} \\
\midrule
قحطی اوکراین (هولودومور ۱۹۳۲-۳۳) & ۳-۷ میلیون \\
پاک‌سازی بزرگ (۱۹۳۶-۳۸) & ۷۰۰ هزار اعدام + میلیون‌ها زندانی \\
گولاگ (کل دوره) & ۱-۱.۵ میلیون مرگ در اردوگاه \\
تبعید ملت‌ها & صدها هزار \\
\midrule
\textbf{جمع تقریبی} & ۶-۱۰ میلیون \\
\bottomrule
\end{tabular}
\end{table}

\begin{warningbox}[یادآوری تاریخی]
ارقام بالا مورد اختلاف مورخان است. برخی ارقام بالاتر (تا ۲۰ میلیون) و برخی پایین‌تر ارائه می‌دهند. اما در اصل فاجعه انسانی دوران استالین اختلافی نیست. این فجایع به نام «کمونیسم» انجام شد، اما بسیاری از کمونیست‌ها (از جمله تروتسکی و دیگران) قربانی آن بودند.
\end{warningbox}

\subsection{تروتسکی در برابر استالین}

\begin{personbox}[لئون تروتسکی (۱۸۷۹-۱۹۴۰)]

\textbf{نام اصلی:} لو داویدوویچ برونشتاین

\textbf{ملیت:} اوکراینی (یهودی)

\textbf{نقش:} رهبر انقلاب اکتبر، بنیان‌گذار ارتش سرخ

\textbf{آثار کلیدی:} \book{انقلاب مداوم} (۱۹۳۰)، \book{انقلاب خیانت‌شده} (۱۹۳۶)

\mediumspace

\textbf{نظریه انقلاب مداوم:}
\begin{itemize}[nosep]
    \item در کشورهای عقب‌مانده، بورژوازی ضعیف است
    \item انقلاب دموکراتیک مستقیماً به انقلاب سوسیالیستی تبدیل می‌شود
    \item اما انقلاب باید \textbf{جهانی} شود، وگرنه منزوی و فاسد می‌شود
\end{itemize}

\mediumspace

\textbf{نقد استالین:} شوروی به «دولت کارگری منحط» تبدیل شده؛ بوروکراسی جای بورژوازی را گرفته.

\mediumspace

\textbf{سرنوشت:} تبعید (۱۹۲۹)، ترور توسط عامل استالین در مکزیک (۱۹۴۰)
\end{personbox}

\begin{comparebox}[مقایسه: استالین و تروتسکی]
\begin{table}[H]
\centering
\begin{tabular}{r|c|c}
\toprule
\textbf{موضوع} & \textbf{استالین} & \textbf{تروتسکی} \\
\midrule
جغرافیای انقلاب & سوسیالیسم در یک کشور & انقلاب مداوم و جهانی \\
صنعتی‌سازی & اجباری و سریع & لازم اما انسانی‌تر \\
بوروکراسی & ابزار ضروری & انحراف از انقلاب \\
دموکراسی حزبی & تمرکز مطلق & لازم (نسبتاً) \\
\bottomrule
\end{tabular}
\end{table}
\end{comparebox}

% ══════════════════════════════════════════════════════════════════════════════
\section{مائوئیسم}
% ══════════════════════════════════════════════════════════════════════════════

\person{مائو تسه‌تونگ} مارکسیسم-لنینیسم را به شرایط چین تطبیق داد و نسخه جدیدی ایجاد کرد که در جهان سوم تأثیرگذار شد.

\begin{personbox}[مائو تسه‌تونگ (۱۸۹۳-۱۹۷۶)]

\textbf{ملیت:} چینی

\textbf{نقش:} رهبر انقلاب چین، رئیس جمهوری خلق چین (۱۹۴۹-۱۹۷۶)

\textbf{آثار کلیدی:} \book{درباره تضاد}، \book{درباره عمل}، \book{کتاب سرخ کوچک}

\mediumspace

\textbf{نوآوری‌های مائو:}
\begin{enumerate}[nosep]
    \item \textbf{انقلاب دهقانی:} دهقانان، نه پرولتاریای شهری، نیروی اصلی انقلاب در کشورهای عقب‌مانده
    \item \textbf{جنگ خلقی طولانی‌مدت:} محاصره شهر از روستا
    \item \textbf{خط توده:} «از توده‌ها به توده‌ها»
    \item \textbf{تضاد اصلی و فرعی:} تشخیص تضاد اصلی در هر مرحله
    \item \textbf{انقلاب فرهنگی:} مبارزه طبقاتی در دوران سوسیالیسم ادامه دارد
\end{enumerate}

\mediumspace

\textbf{نقل‌قول:} \textit{«قدرت سیاسی از لوله تفنگ بیرون می‌آید.»}
\end{personbox}

\begin{figure}[htbp]
\centering
\begin{tikzpicture}[scale=0.85]

% عنوان
\node[font=\large\bfseries, text=DarkGray] at (0, 7) {\fa{مقایسه: لنینیسم و مائوئیسم}};

% کادر لنینیسم
\fill[LeftRed, opacity=0.15, rounded corners=10pt] (-7, 0.5) rectangle (-0.5, 6);
\draw[line width=2pt, color=LeftRed, rounded corners=10pt] (-7, 0.5) rectangle (-0.5, 6);

\node[text=LeftRed, font=\bfseries\large] at (-3.75, 5.5) {\fa{لنینیسم}};

\node[text=DarkGray, font=\small, text width=5.5cm, align=center] at (-3.75, 4) {
    \fa{پرولتاریای صنعتی = نیروی انقلاب}\\[5pt]
    \fa{حزب پیشاهنگ رهبری می‌کند}\\[5pt]
    \fa{قیام شهری}\\[5pt]
    \fa{کشورهای صنعتی اولویت دارند}\\[5pt]
    \fa{انترناسیونالیسم پرولتری}
};

\node[text=LeftRedDark, font=\small\bfseries] at (-3.75, 1.2) {\fa{زمینه: روسیه ۱۹۱۷}};

% کادر مائوئیسم
\fill[OrangeAccent, opacity=0.15, rounded corners=10pt] (0.5, 0.5) rectangle (7, 6);
\draw[line width=2pt, color=OrangeAccent, rounded corners=10pt] (0.5, 0.5) rectangle (7, 6);

\node[text=OrangeAccent, font=\bfseries\large] at (3.75, 5.5) {\fa{مائوئیسم}};

\node[text=DarkGray, font=\small, text width=5.5cm, align=center] at (3.75, 4) {
    \fa{دهقانان = نیروی اصلی انقلاب}\\[5pt]
    \fa{حزب در خدمت توده‌ها}\\[5pt]
    \fa{جنگ چریکی طولانی}\\[5pt]
    \fa{جهان سوم مرکز انقلاب}\\[5pt]
    \fa{ضد امپریالیسم ملی}
};

\node[text=OrangeAccent!80!black, font=\small\bfseries] at (3.75, 1.2) {\fa{زمینه: چین ۱۹۲۷-۱۹۴۹}};

\end{tikzpicture}
\caption{تفاوت‌های لنینیسم و مائوئیسم}
\label{fig:leninism-maoism}
\end{figure}

\subsection{انقلاب فرهنگی (۱۹۶۶-۱۹۷۶)}

\begin{historybox}[انقلاب فرهنگی پرولتاریایی بزرگ]
مائو در ۱۹۶۶، پس از شکست «جهش بزرگ به پیش»، انقلاب فرهنگی را آغاز کرد:
\begin{itemize}[nosep]
    \item هدف رسمی: مبارزه با «راه سرمایه‌داری» در حزب
    \item \textbf{گاردهای سرخ:} جوانان بسیج شدند
    \item روشنفکران، مقامات، معلمان مورد حمله قرار گرفتند
    \item میراث فرهنگی نابود شد
    \item هرج‌ومرج، خشونت، میلیون‌ها قربانی
\end{itemize}
\textbf{نتیجه:} ویرانی اقتصادی و فرهنگی. پس از مرگ مائو (۱۹۷۶)، «دار و دسته چهار نفره» محاکمه شدند و چین مسیر متفاوتی در پیش گرفت.
\end{historybox}

% ══════════════════════════════════════════════════════════════════════════════
\section{کمونیسم در جهان سوم}
% ══════════════════════════════════════════════════════════════════════════════

کمونیسم در قرن بیستم بیش از اروپا در جهان سوم موفق شد—جایی که مارکس کمتر انتظارش را داشت.

\subsection{انقلاب کوبا}

\begin{personbox}[فیدل کاسترو (۱۹۲۶-۲۰۱۶)]

\textbf{ملیت:} کوبایی

\textbf{نقش:} رهبر انقلاب کوبا، حکومت ۱۹۵۹-۲۰۰۸

\mediumspace

\textbf{ویژگی‌های انقلاب کوبا:}
\begin{itemize}[nosep]
    \item انقلاب ابتدا «کمونیستی» نبود—ملی‌گرا و ضد دیکتاتوری بود
    \item فشار آمریکا کوبا را به سوی شوروی راند
    \item بحران موشکی ۱۹۶۲
    \item مدل: بهداشت و آموزش رایگان، اما فقدان آزادی سیاسی
\end{itemize}
\end{personbox}

\begin{personbox}[ارنستو چه گوارا (۱۹۲۸-۱۹۶۷)]

\textbf{ملیت:} آرژانتینی

\textbf{نقش:} انقلابی، نظریه‌پرداز جنگ چریکی

\textbf{اثر کلیدی:} \book{جنگ چریکی} (۱۹۶۱)

\mediumspace

\textbf{ایده‌های چه:}
\begin{itemize}[nosep]
    \item \textbf{فوکوئیسم:} گروه کوچک مسلح می‌تواند «کانون» انقلاب شود
    \item انترناسیونالیسم: «دو، سه، بسیار ویتنام بسازیم»
    \item نقد بوروکراسی شوروی
    \item «انسان نوین سوسیالیستی»
\end{itemize}

\mediumspace

\textbf{سرنوشت:} کشته شده توسط ارتش بولیوی با کمک سیا (۱۹۶۷). نماد جنبش چپ جهانی شد.
\end{personbox}

\subsection{ویتنام}

\begin{historybox}[جنگ ویتنام و پیروزی کمونیسم]
\begin{itemize}[nosep]
    \item \person{هو شی‌مین}: رهبر ویت‌مین، ترکیب ناسیونالیسم و کمونیسم
    \item ۱۹۵۴: شکست فرانسه در دین‌بین‌فو
    \item ۱۹۵۵-۱۹۷۵: جنگ با آمریکا
    \item ۱۹۷۵: سقوط سایگون، وحدت ویتنام تحت حکومت کمونیست
\end{itemize}
\textbf{اهمیت:} اولین شکست نظامی آمریکا؛ الهام‌بخش جنبش‌های ضد امپریالیستی جهان
\end{historybox}

% ══════════════════════════════════════════════════════════════════════════════
\section{یوروکمونیسم}
% ══════════════════════════════════════════════════════════════════════════════

در اروپای غربی، برخی احزاب کمونیست کوشیدند از الگوی شوروی فاصله بگیرند.

\begin{defbox}[تعریف: یوروکمونیسم]
گرایشی در احزاب کمونیست اروپای غربی (ایتالیا، فرانسه، اسپانیا) در دهه ۱۹۷۰:
\begin{itemize}[nosep]
    \item پذیرش دموکراسی پارلمانی و کثرت‌گرایی
    \item استقلال از مسکو
    \item رد دیکتاتوری پرولتاریا
    \item «راه ملی به سوسیالیسم»
\end{itemize}
\textbf{چهره‌های کلیدی:} انریکو برلینگوئر (ایتالیا)، سانتیاگو کاریو (اسپانیا)
\end{defbox}

\begin{personbox}[آنتونیو گرامشی (۱۸۹۱-۱۹۳۷)]

\textbf{ملیت:} ایتالیایی

\textbf{اثر کلیدی:} \book{دفترهای زندان}

\textbf{ایده محوری:} هژمونی فرهنگی

\mediumspace

\person{گرامشی}، رهبر حزب کمونیست ایتالیا، در زندان فاشیستی نظریه‌ای نو ارائه داد:
\begin{itemize}[nosep]
    \item \textbf{هژمونی:} طبقه حاکم نه فقط با زور، بلکه با «رضایت» حکومت می‌کند
    \item این رضایت از طریق فرهنگ، رسانه، آموزش ساخته می‌شود
    \item انقلاب نیاز به \textbf{جنگ موضعی} فرهنگی دارد، نه فقط حمله به دولت
    \item \textbf{روشنفکر ارگانیک:} روشنفکر متعهد به طبقه
\end{itemize}

\mediumspace

\textbf{تأثیر:} چپ نو، مطالعات فرهنگی، یوروکمونیسم همگی از گرامشی الهام گرفتند.
\end{personbox}

% ══════════════════════════════════════════════════════════════════════════════
\section{فروپاشی بلوک شرق}
% ══════════════════════════════════════════════════════════════════════════════

\begin{figure}[htbp]
\centering
\begin{tikzpicture}[scale=0.85]

% عنوان
\node[font=\large\bfseries, text=DarkGray] at (0, 7.5) {\fa{تایم‌لاین: فروپاشی کمونیسم در اروپا (۱۹۸۹-۱۹۹۱)}};

% خط زمان
\draw[line width=4pt, MediumGray] (-7, 5.5) -- (7, 5.5);

% نقاط
% لهستان
\fill[CenterGreen] (-6, 5.5) circle (9pt);
\node[above, text=CenterGreen, font=\scriptsize\bfseries] at (-6, 6) {\fa{ژوئن ۸۹}};
\node[below, text width=2cm, align=center, text=DarkGray, font=\tiny] at (-6, 4.9) {
    \fa{لهستان}\\
    \fa{انتخابات آزاد}
};

% مجارستان
\fill[CenterGreen] (-4, 5.5) circle (9pt);
\node[above, text=CenterGreen, font=\scriptsize\bfseries] at (-4, 6) {\fa{اکتبر ۸۹}};
\node[below, text width=2cm, align=center, text=DarkGray, font=\tiny] at (-4, 4.9) {
    \fa{مجارستان}\\
    \fa{جمهوری}
};

% دیوار برلین
\fill[GoldAccent] (-1.5, 5.5) circle (11pt);
\node[above, text=GoldAccent, font=\scriptsize\bfseries] at (-1.5, 6.1) {\fa{۹ نوامبر ۸۹}};
\node[below, text width=2.2cm, align=center, text=DarkGray, font=\tiny] at (-1.5, 4.8) {
    \fa{سقوط}\\
    \fa{دیوار برلین}
};

% چکسلواکی
\fill[CenterGreen] (1, 5.5) circle (9pt);
\node[above, text=CenterGreen, font=\scriptsize\bfseries] at (1, 6) {\fa{نوامبر ۸۹}};
\node[below, text width=2cm, align=center, text=DarkGray, font=\tiny] at (1, 4.9) {
    \fa{چکسلواکی}\\
    \fa{انقلاب مخملی}
};

% رومانی
\fill[LeftRed] (3, 5.5) circle (9pt);
\node[above, text=LeftRed, font=\scriptsize\bfseries] at (3, 6) {\fa{دسامبر ۸۹}};
\node[below, text width=2cm, align=center, text=DarkGray, font=\tiny] at (3, 4.9) {
    \fa{رومانی}\\
    \fa{چائوشسکو اعدام}
};

% آلمان
\fill[GoldAccent] (5, 5.5) circle (9pt);
\node[above, text=GoldAccent, font=\scriptsize\bfseries] at (5, 6) {\fa{اکتبر ۹۰}};
\node[below, text width=2cm, align=center, text=DarkGray, font=\tiny] at (5, 4.9) {
    \fa{وحدت}\\
    \fa{آلمان}
};

% شوروی
\fill[LeftRedDark] (7, 5.5) circle (11pt);
\node[above, text=LeftRedDark, font=\scriptsize\bfseries] at (7, 6.1) {\fa{دسامبر ۹۱}};
\node[below, text width=2cm, align=center, text=DarkGray, font=\tiny] at (7, 4.8) {
    \fa{فروپاشی}\\
    \fa{شوروی}
};

% توضیح
\node[text=MediumGray, font=\small] at (0, 3.5) {\fa{در کمتر از سه سال، کل بلوک شرق فروپاشید.}};

\end{tikzpicture}
\caption{فروپاشی کمونیسم در اروپا}
\label{fig:communism-collapse}
\end{figure}

\subsection{علل فروپاشی}

\begin{debatebox}[چرا شوروی فروپاشید؟]
\textbf{تبیین‌های اقتصادی:}
\begin{itemize}[nosep]
    \item ناکارآمدی برنامه‌ریزی مرکزی
    \item رقابت تسلیحاتی با آمریکا
    \item افت قیمت نفت در دهه ۱۹۸۰
\end{itemize}

\textbf{تبیین‌های سیاسی:}
\begin{itemize}[nosep]
    \item فقدان مشروعیت دموکراتیک
    \item اصلاحات گورباچف (گلاسنوست، پرسترویکا) کنترل را از بین برد
    \item ناسیونالیسم ملت‌های غیرروس
\end{itemize}

\textbf{تبیین‌های ایدئولوژیک:}
\begin{itemize}[nosep]
    \item ایمان به کمونیسم از بین رفته بود
    \item نخبگان خودشان می‌خواستند سرمایه‌دار شوند
\end{itemize}
\end{debatebox}

\subsection{ارزیابی‌های متفاوت از فروپاشی}

\begin{comparebox}[فروپاشی شوروی: پایان چه چیزی؟]
\begin{table}[H]
\centering
\small
\begin{tabular}{r|c|c}
\toprule
\textbf{دیدگاه} & \textbf{ارزیابی} & \textbf{نماینده} \\
\midrule
لیبرال & پایان تاریخ، پیروزی دموکراسی لیبرال & فوکویاما \\
\midrule
محافظه‌کار & پیروزی غرب در جنگ سرد & ریگان، تاچر \\
\midrule
مارکسیست منتقد & شکست استالینیسم، نه مارکسیسم & تروتسکیست‌ها \\
\midrule
چپ نو & فرصت بازاندیشی در سوسیالیسم & — \\
\bottomrule
\end{tabular}
\end{table}
\end{comparebox}

% ══════════════════════════════════════════════════════════════════════════════
\section{کمونیسم امروز}
% ══════════════════════════════════════════════════════════════════════════════

\subsection{چین: کمونیسم یا سرمایه‌داری دولتی؟}

\begin{debatebox}[چین امروز چیست؟]
\textbf{دفاع حزب کمونیست چین:}
\begin{itemize}[nosep]
    \item «سوسیالیسم با ویژگی‌های چینی»
    \item اقتصاد بازار ابزار است، نه هدف
    \item هنوز در مرحله اولیه سوسیالیسم هستیم
    \item حزب کمونیست رهبری می‌کند
\end{itemize}

\textbf{نقد از چپ:}
\begin{itemize}[nosep]
    \item این سرمایه‌داری دولتی است، نه سوسیالیسم
    \item استثمار کارگران شدید است
    \item نابرابری به شدت افزایش یافته
    \item کجاست کنترل کارگران بر تولید؟
\end{itemize}

\textbf{نقد از راست:}
\begin{itemize}[nosep]
    \item این سرمایه‌داری اقتدارگراست
    \item موفقیت چین ناشی از بازار است، نه کمونیسم
\end{itemize}
\end{debatebox}

\subsection{کوبا، ویتنام، کره شمالی}

\begin{table}[htbp]
\centering
\caption{کشورهای کمونیست امروز}
\label{tab:communist-countries-today}
\begin{tabular}{r|c|c|c}
\toprule
\textbf{کشور} & \textbf{حزب حاکم} & \textbf{سیستم اقتصادی} & \textbf{ویژگی} \\
\midrule
چین & کمونیست & اقتصاد بازار & ابرقدرت اقتصادی \\
ویتنام & کمونیست & اقتصاد بازار (دوی‌موی) & رشد سریع \\
کوبا & کمونیست & مختلط (اصلاحات) & تحت تحریم آمریکا \\
لائوس & انقلابی خلق & اقتصاد بازار & کم‌جمعیت \\
کره شمالی & کارگران (جوچه) & برنامه‌ریزی مرکزی & منزوی، فقیر \\
\bottomrule
\end{tabular}
\end{table}

% ══════════════════════════════════════════════════════════════════════════════
%                         نقشه کشورهای کمونیستی
% ══════════════════════════════════════════════════════════════════════════════

\begin{landscape}
\begin{figure}[htbp]
\centering
\begin{tikzpicture}[scale=1.2]

% عنوان
\node[font=\Large\bfseries, text=DarkGray] at (8, 7) {\fa{کشورهای کمونیستی در تاریخ}};

% راهنما
\fill[LeftRedDark, opacity=0.8] (0, -0.5) rectangle (1, 0);
\node[left, text=DarkGray, font=\small] at (2.8, -0.25) {\fa{هنوز کمونیست}};

\fill[LeftRed, opacity=0.5] (4, -0.5) rectangle (5, 0);
\node[left, text=DarkGray, font=\small] at (6.8, -0.25) {\fa{سابقاً کمونیست}};

\fill[OrangeAccent, opacity=0.5] (8, -0.5) rectangle (9, 0);
\node[left, text=DarkGray, font=\small] at (11.5, -0.25) {\fa{حکومت کمونیستی موقت}};

% ────────── اروپا ──────────
\node[font=\bfseries, text=RightBlue] at (2, 5.5) {\fa{اروپا}};

% شوروی
\fill[LeftRed, opacity=0.5, rounded corners=3pt] (1, 4) rectangle (5, 5);
\node[text=DarkGray, font=\scriptsize] at (3, 4.5) {\fa{شوروی (۱۹۲۲-۱۹۹۱)}};

% اروپای شرقی
\fill[LeftRed, opacity=0.4, rounded corners=3pt] (0.5, 2.5) rectangle (4, 3.5);
\node[text=DarkGray, font=\tiny, text width=3cm, align=center] at (2.25, 3) {
    \fa{لهستان، چکسلواکی، مجارستان}\\
    \fa{آلمان شرقی، رومانی، بلغارستان}\\
    \fa{(۱۹۴۵-۱۹۸۹)}
};

% یوگسلاوی
\fill[OrangeAccent, opacity=0.4, rounded corners=3pt] (0.5, 1.5) rectangle (3, 2.2);
\node[text=DarkGray, font=\tiny] at (1.75, 1.85) {\fa{یوگسلاوی (۱۹۴۵-۱۹۹۲)}};

% آلبانی
\fill[LeftRed, opacity=0.4, rounded corners=3pt] (3.2, 1.5) rectangle (5, 2.2);
\node[text=DarkGray, font=\tiny] at (4.1, 1.85) {\fa{آلبانی (۱۹۴۴-۱۹۹۱)}};

% ────────── آسیا ──────────
\node[font=\bfseries, text=RightBlue] at (9, 5.5) {\fa{آسیا}};

% چین
\fill[LeftRedDark, opacity=0.8, rounded corners=3pt] (7, 3.5) rectangle (11, 5);
\node[text=white, font=\small\bfseries] at (9, 4.25) {\fa{چین (۱۹۴۹-اکنون)}};

% ویتنام
\fill[LeftRedDark, opacity=0.8, rounded corners=3pt] (7, 2.5) rectangle (9.5, 3.2);
\node[text=white, font=\scriptsize] at (8.25, 2.85) {\fa{ویتنام (۱۹۷۵-اکنون)}};

% کره شمالی
\fill[LeftRedDark, opacity=0.8, rounded corners=3pt] (10, 2.5) rectangle (12.5, 3.2);
\node[text=white, font=\scriptsize] at (11.25, 2.85) {\fa{کره شمالی (۱۹۴۸-اکنون)}};

% لائوس
\fill[LeftRedDark, opacity=0.7, rounded corners=3pt] (7, 1.7) rectangle (9, 2.3);
\node[text=white, font=\tiny] at (8, 2) {\fa{لائوس}};

% کامبوج
\fill[LeftRed, opacity=0.4, rounded corners=3pt] (9.2, 1.7) rectangle (11.5, 2.3);
\node[text=DarkGray, font=\tiny] at (10.35, 2) {\fa{کامبوج (۷۵-۷۹)}};

% افغانستان
\fill[LeftRed, opacity=0.4, rounded corners=3pt] (11.7, 1.7) rectangle (14, 2.3);
\node[text=DarkGray, font=\tiny] at (12.85, 2) {\fa{افغانستان (۷۸-۹۲)}};

% ────────── آمریکا ──────────
\node[font=\bfseries, text=RightBlue] at (14, 5.5) {\fa{آمریکا}};

% کوبا
\fill[LeftRedDark, opacity=0.8, rounded corners=3pt] (13, 4) rectangle (15.5, 5);
\node[text=white, font=\small\bfseries] at (14.25, 4.5) {\fa{کوبا (۱۹۵۹-اکنون)}};

% نیکاراگوئه
\fill[OrangeAccent, opacity=0.5, rounded corners=3pt] (13, 3) rectangle (15.5, 3.7);
\node[text=DarkGray, font=\tiny] at (14.25, 3.35) {\fa{نیکاراگوئه (۱۹۷۹-۱۹۹۰)}};

% ────────── آفریقا ──────────
\node[font=\bfseries, text=RightBlue] at (6, 0.5) {\fa{آفریقا}};

\fill[OrangeAccent, opacity=0.4, rounded corners=3pt] (4.5, -0.8) rectangle (8, 0);
\node[text=DarkGray, font=\tiny, text width=3cm, align=center] at (6.25, -0.4) {
    \fa{اتیوپی، موزامبیک، آنگولا}\\
    \fa{(دهه‌های ۷۰-۸۰)}
};

\end{tikzpicture}
\caption{گستره جغرافیایی کمونیسم در قرن بیستم}
\label{fig:communist-map}
\end{figure}
\end{landscape}

% ══════════════════════════════════════════════════════════════════════════════
%                         خلاصه و پرسش‌ها
% ══════════════════════════════════════════════════════════════════════════════

\newpage

\begin{summarybox}

\textbf{۱. تمایز مفهومی:} سوسیالیسم مرحله گذار است؛ کمونیسم جامعه نهایی بی‌طبقه و بی‌دولت.

\textbf{۲. لنینیسم:} حزب پیشاهنگ، امپریالیسم، انقلاب در «ضعیف‌ترین حلقه». انقلاب اکتبر ۱۹۱۷ مارکسیسم را به قدرت رساند.

\textbf{۳. استالینیسم:} «سوسیالیسم در یک کشور»، صنعتی‌سازی اجباری، پاک‌سازی‌ها. میلیون‌ها قربانی.

\textbf{۴. تروتسکیسم:} انقلاب مداوم، نقد بوروکراسی شوروی، انترناسیونال چهارم.

\textbf{۵. مائوئیسم:} انقلاب دهقانی، جنگ خلقی، انقلاب فرهنگی. تأثیرگذار در جهان سوم.

\textbf{۶. گسترش جهانی:} کوبا، ویتنام، آفریقا. کمونیسم در جهان سوم موفق‌تر از اروپا بود.

\textbf{۷. یوروکمونیسم:} تلاش برای ترکیب کمونیسم با دموکراسی. گرامشی و هژمونی فرهنگی.

\textbf{۸. فروپاشی:} ۱۹۸۹-۱۹۹۱ بلوک شرق فروپاشید. علل: اقتصادی، سیاسی، ایدئولوژیک.

\textbf{۹. امروز:} چین، ویتنام، کوبا، لائوس، کره شمالی. اما چین و ویتنام عملاً اقتصاد بازار دارند.

\end{summarybox}

\vspace{20pt}

\begin{questionbox}

\begin{enumerate}[nosep, rightmargin=1cm]
    \item آیا استالینیسم نتیجه منطقی لنینیسم بود، یا انحرافی از آن؟ آیا می‌شد «لنینیسم بدون استالینیسم» داشت؟

    \item تروتسکی گفت انقلاب باید جهانی شود وگرنه منحط می‌شود. آیا تجربه شوروی او را تأیید کرد؟

    \item چین امروز را چه باید نامید: «سوسیالیسم با ویژگی‌های چینی»، «سرمایه‌داری دولتی»، یا چیز دیگر؟

    \item با توجه به تجربه قرن بیستم، آیا «کمونیسم» (جامعه بی‌طبقه و بی‌دولت) هنوز هدفی معقول است؟
\end{enumerate}

\end{questionbox}

% ══════════════════════════════════════════════════════════════════════════════
%                         پایان فصل پنجم
% ══════════════════════════════════════════════════════════════════════════════